% =====================================================================
% Section 3: Model and Mechanism
%
% Structure:
%   3.1  Physical motivation (origin story, honest)
%   3.2  Network model (graph, attributes, potential function)
%   3.3  Per-packet state and mass dynamics
%   3.4  Thermal memory and global thermostat
% =====================================================================

\section{Model and Mechanism}
\label{sec:model}

This section presents the mathematical model underlying mass-aware
routing. We begin with the physical intuition that motivated its
form (\S\ref{subsec:physical-motivation}), then formalize the network
model (\S\ref{subsec:network-model}), the per-packet state dynamics
(\S\ref{subsec:mass-dynamics}), and the thermal memory mechanism
(\S\ref{subsec:thermal-memory}).

% =====================================================================
\subsection{Physical motivation}
\label{subsec:physical-motivation}

The functional form of the algorithm presented in this paper
originated as a thought experiment about the physical analog of a
packet in transit. We record the analogy here because it shaped the
mathematical decisions that follow---the multiplicative mass update,
the integrated potential, the role of a global temperature---and
because we believe the analogy retains explanatory value beyond its
historical role. We emphasize that the algorithm's correctness and
empirical performance, established in
Sections~\ref{sec:algorithm}--\ref{sec:results}, do not rely on the
analogy: the physical motivation is heuristic, not foundational. In
particular, the scalar we call \emph{mass} below is closer in spirit
to the \emph{effective mass} of solid-state physics---a quantity that
varies with the medium a particle traverses---than to the invariant
inertial mass of classical mechanics or to any thermodynamic quantity.
The intuition is the one motivating the original effective-mass
picture: when a metal is heated, the crystal lattice expands and the
conduction electrons interact differently with that perturbed lattice;
solid-state physicists model the change as a varying $m^{*}$ rather
than as a change in the electron's true mass. We borrow only the
structural form of that idea: a scalar that grows with the medium's
local conditions and modifies the cost of further interaction.

Consider a single packet traversing a sequence of network links.
Treating the packet as a particle moving through a medium suggests
several mappings, each of which has a counterpart in routing:

\begin{itemize}
\item \emph{Velocity and friction.} A particle moving through a
denser medium dissipates more energy per unit distance. In the
network analog, a packet traversing a more loaded edge encounters
more ``friction''---a higher probability of queueing, drop, or
retransmission. The local congestion ratio $\rho(u,v) = \lambda(u,v)/c(u,v)$
plays the role of a medium density.

\item \emph{Heat accumulation.} A particle that experiences friction
warms the medium it traverses, leaving a thermal trace. In the
network analog, a packet's transit through a congested edge raises
the probability that subsequent packets crossing the same edge
will experience loss. This thermal trace---the per-edge loss
history---we encode as the snapshot $\theta(u,v)$, decayed
exponentially over time.

\item \emph{Effective mass.} A particle whose trajectory has crossed turbulent regions
accumulates an effective resistance to further perturbation; in
our analogy we encode this as a growing ``effective mass'' that
makes the packet costlier to push through additional congestion.
Translating this to routing: a packet
that has already traversed congested edges has, in some sense,
done more work against the network, and any further passage
through congested regions should be discouraged. We encode this
accumulated history as a scalar $m\ge 1$ that the packet carries
internally and that grows multiplicatively whenever the packet
traverses a congested edge:
\begin{equation}
m_{\text{out}} = \min\!\Big(m_{\text{in}}\cdot
                       \big(1+\kappa\,\rho(u,v)\big),\;m_{\max}\Big).
\label{eq:mass-update}
\end{equation}

\item \emph{Global temperature and thermostat coupling.} A
hot system accelerates locally hot processes; the network analog
is that when overall load is high, congestion penalties should
weigh more heavily on routing decisions even at edges that are
not themselves the most loaded. We capture this by defining a
global temperature $T(G) = |E|^{-1}\sum_{(u,v)\in E}\rho(u,v)$
and coupling the congestion coefficient through
$\beta_{\text{eff}} = \beta(1 + \theta_T\cdot T(G))$.
\end{itemize}

Three structural commitments emerged from this analogy and are
reflected throughout the paper:

\begin{enumerate}
\item \textbf{Per-packet state.} Each packet carries information
about its own trajectory. This is the central commitment: routing
decisions for a packet depend not only on the network state at the
moment of decision, but also on the cumulative experience of the
packet itself.

\item \textbf{Multiplicative degradation.} The mass update is
multiplicative rather than additive
(Eq.~\ref{eq:mass-update}). This choice reflects the physical
intuition that the cost of further congestion compounds with prior
congestion: a fresh packet is robust to the first congested edge,
but a heavily-laden packet should be steered around even mildly
congested neighborhoods.

\item \textbf{Saturation.} The mass is capped at $m_{\max}$. Beyond
saturation, additional congestion does not change the packet's
routing pressure; this avoids pathological behavior where one
extremely degraded packet would dominate routing decisions for an
arbitrarily long path.
\end{enumerate}

These three commitments---per-packet state, multiplicative growth,
and saturation---are the operational content of the physical
analogy. The remainder of this section makes them precise.

% =====================================================================
\subsection{Network model}
\label{subsec:network-model}

We model a network as an undirected graph $G=(V,E)$ with $|V|=n$
nodes and $|E|=e$ edges. Each edge $(u,v)\in E$ carries four
time-varying attributes:
\begin{itemize}
\item \emph{Latency} $\ell(u,v) > 0$: a fixed propagation cost
sampled per realization from a positive distribution.
\item \emph{Capacity} $c(u,v) > 0$: maximum simultaneous load
the edge can sustain before drop.
\item \emph{Load} $\lambda(u,v) \ge 0$: instantaneous load on
the edge, updated as packets traverse and decayed by a factor
$d_\lambda \in (0,1)$ at each routing step.
\item \emph{Loss snapshot} $\theta(u,v) \ge 0$: cumulative
loss memory, decayed by a factor $d_\theta \in (0,1)$ each step.
\end{itemize}

We define the local congestion of edge $(u,v)$ as
$\rho(u,v) = \lambda(u,v)/c(u,v) \in [0,1]$ when the edge is
operating below capacity. The global temperature of the network
is the mean local congestion:
\begin{equation}
T(G) = \frac{1}{|E|}\sum_{(u,v)\in E}\rho(u,v).
\label{eq:global-temperature}
\end{equation}

The \emph{integrated edge potential} combines latency, congestion,
and thermal memory into a single scalar cost:
\begin{equation}
\Phi(u,v) = \alpha\,\ell(u,v)
            + \beta_{\text{eff}}\,\rho(u,v)
            + \gamma\,\theta(u,v),
\label{eq:edge-potential}
\end{equation}
where $\beta_{\text{eff}} = \beta\,(1 + \theta_T\,T(G))$ is the
thermostat-coupled congestion weight, and
$\alpha,\beta,\gamma,\theta_T$ are non-negative coefficients
controlling the relative contribution of each term. Throughout
this paper we use $\alpha=1$, $\beta=3$, $\gamma=2$,
$\theta_T=5$, $d_\lambda = 0.9$, $d_\theta=0.999$
(see \S\ref{sec:experiments} for justification of these defaults).

% =====================================================================
\subsection{Per-packet state and mass dynamics}
\label{subsec:mass-dynamics}

A central novelty of this work is that each packet carries an
internal scalar state, the \emph{mass} $m$, alongside its source
and destination. The mass is initialized at $m_0 = 1$ when the
packet is issued and evolves multiplicatively as the packet
traverses edges, according to Equation~\ref{eq:mass-update}. The
parameter $\kappa\ge 0$ controls the rate of growth; $\kappa = 0$
disables the mechanism, reducing the algorithm to its stateless
counterpart, while $\kappa \to \infty$ saturates the mass after a
single congested edge. Empirically we find
$\kappa \in [0.5, 1.5]$ is the productive range, with $\kappa=1$
being Pareto-optimal across the topologies we evaluate
(\S\ref{sec:results}).

The mass enters the routing decision through a multiplicative
scaling of the edge potential:
\begin{equation}
\delta(u,v\mid m) = m\cdot \Phi(u,v).
\label{eq:cost-step}
\end{equation}
A packet with mass $m=1.0$ pays the bare potential; a packet with
mass $m=2.0$ pays double for every edge it considers. The
consequence is that a packet that has accumulated mass through
congested edges is, by construction, more expensive to route
through any further congested edge, and the dynamic-programming
router (Algorithm~\ref{alg:massaware-dp}) systematically prefers
cooler alternatives for it.

The mass is bounded above by $m_{\max}$, ensuring that the
multiplicative growth cannot cause numerical or behavioral
pathologies. We use $m_{\max}=3$ throughout this paper, a value
empirically sufficient to discriminate routing decisions while
keeping the bucket discretization (Algorithm~\ref{alg:massaware-dp})
tractable.

For dynamic programming purposes, the continuous mass scalar is
discretized into $B$ buckets uniformly distributed in $[1, m_{\max}]$;
all experiments in this paper use $B=32$, with sensitivity to $B$
analyzed in \S\ref{sec:results}.

% =====================================================================
\subsection{Thermal memory and global thermostat}
\label{subsec:thermal-memory}

Two mechanisms in our model carry information about past network
behavior into future routing decisions: the per-edge loss snapshot
$\theta(u,v)$, and the global temperature $T(G)$.

The loss snapshot tracks the cumulative drop history at each edge,
incremented whenever a packet attempts to traverse the edge with
$\lambda(u,v)\ge c(u,v)$ and decayed by a factor $d_\theta$ at each
simulation step. Because the snapshot decays exponentially, it
captures recent loss events more strongly than ancient ones,
acting as an exponentially-weighted moving average of edge
unreliability. The coefficient $\gamma$ in
Equation~\ref{eq:edge-potential} controls how strongly recent
losses bias the routing potential.

The global temperature $T(G)$ aggregates per-edge congestion into
a scalar that reflects the overall stress on the network. When
$T(G)$ is high, the thermostat-coupled coefficient
$\beta_{\text{eff}}$ amplifies the per-edge congestion penalty,
shifting the algorithm into a more cautious regime that prefers
underutilized edges even at the cost of additional latency. When
$T(G)$ is low, $\beta_{\text{eff}}\to\beta$ and the algorithm
behaves as a standard load-aware router.

Together, the loss snapshot and the global thermostat give the
algorithm two distinct memory horizons: a slow, edge-local memory
that tracks chronic unreliability, and a fast, network-global
indicator that responds within a few steps to bursts of
congestion. The integrated potential $\Phi(u,v)$ in
Equation~\ref{eq:edge-potential} combines both, and the
mass dynamics in Equation~\ref{eq:mass-update} weights them by the
packet's accumulated trajectory cost.

This three-way coupling---per-edge snapshot, global thermostat,
per-packet mass---is what distinguishes mass-aware routing from
prior load-aware schemes and what produces the empirical
improvements documented in \S\ref{sec:results}.

% =====================================================================
\subsection{From v1 to EMMET-Newt: activating the loss snapshot}
\label{subsec:newt-extension}

The model just described stores a persistent per-edge loss snapshot
$\sigma(u,v)$ and decays it at each step, but in the v1 algorithm
this snapshot is populated only during warmup; once the production
workload starts, $\sigma$ decays without ever being incremented.
Under a stationary workload (Poisson) this is acceptable: the
distribution that produced the warmup snapshot matches the
distribution at test time. Under a non-stationary workload (bursty
traffic, time-varying demand) the warmup snapshot becomes stale,
and the $\gamma$ coefficient that biases the integrated potential
operates on a frozen picture of past losses rather than a live one.

We close this gap with what we call the \emph{reaction-memory}
mechanism: every time a packet drops on edge $e$ at simulation
step $t$, we increment $\sigma(e) \leftarrow \sigma(e) +
\rho_{\mathrm{blood}}$ before the routine decay takes effect.
The increment $\rho_{\mathrm{blood}}$ (we use values in $[5, 10]$
in practice) acts as a discrete Newton-III reaction: the action (a
packet dying on an edge) produces an opposed reaction (a persistent
disincentive for subsequent packets to route through that edge).
The mechanism is local to the dying edge and decays at the same
rate as the warmup snapshot, so it integrates into the existing
edge potential without changing the search structure of the DP
router.

When $\rho_{\mathrm{blood}} = 0$, the reaction-memory mechanism
degenerates to the v1 algorithm bit-for-bit: $\sigma(e)$ is not
updated at runtime and the warmup snapshot decays as before.
We have verified this degeneration empirically on every topology
tested. The mechanism is strictly additive: any v1 guarantee
transfers to EMMET-Newt with $\rho_{\mathrm{blood}} = 0$.

Under realistic operating conditions feedback is not instantaneous;
intermediate routers do not learn of a drop the moment it happens.
We therefore evaluate two variants. The \emph{immediate} variant
updates $\sigma$ at the instant a packet dies and gives an upper
bound on the mechanism's potential. The \emph{delayed} variant
buffers drop events during a burst and releases them to $\sigma$ at
the burst boundary, the conservative reading of what is achievable
in practice. The empirical evaluation in \S\ref{sec:results}
reports both, and treats the delayed variant as the headline.
